% !PRINTED 658
\weberpage{658}

Then from (20) a congruence would follow:

\[ y_0 + y_1\tau + y_2\tau^2 + \cdots + y_{n-1}\tau^{n-1} \equiv 0 \quad (\mathrm{mod}\ z-c), \]

which by Theorem 5. is not possible.

By § 174, (17) and § 176, (9) we have

\weq{27}{\begin{aligned}
N F'(\tau) &= \pm \Delta(1, \tau, \tau^2, \ldots, \tau^{n-1}) \\
&= \pm U^2 \Delta(\omega_1, \omega_2, \ldots, \omega_n) \\
&= \pm U^2 \Delta;
\end{aligned}}

thus

\weq{28}{N_a F'(\tau) = \Delta}

is \gesperrt{the field discriminant, which is accordingly decomposed by (24) into
linear factors.}

We use the results obtained for the proof of the following theorem as
well:

7. \gesperrt{If an integral function in $\Omega$}

\weq{29}{\omega = x_1\omega_1 + x_2\omega_2 + \cdots + x_n\omega_n}

\gesperrt{is divisible by every prime functional that divides $z-c$, then
$S(\omega)$ is divisible by $z-c$.}

Here $x_1, x_2, \ldots, x_n$ are integral functions in $Z$. The function
$\omega$ arises from $\tau$ through the substitution

\weq{30}{t_1 = x_1,\ t_2 = x_2,\ \ldots,\ t_n = x_n.}

If therefore $\omega$ is divisible by $\pi$, and $\omega_1 \equiv c_1,
\ldots, \omega_n \equiv c_n \pmod{\pi}$, then

\[ x_1c_1 + x_2c_2 + \cdots + x_nc_n \equiv 0 \quad (\mathrm{mod}\ \pi) \]

and consequently it is also, as a function in $Z$, divisible by $z-c$. If
therefore $\tau_1, \tau_2, \ldots, \tau_m$ are the constant functionals
to which $\tau$ is congruent with respect to the prime factors $\pi_1,
\pi_2, \ldots, \pi_m$ of $(z-c)$, and if these pass by the substitution
(30) into $\tau_1^0, \tau_2^0, \ldots, \tau_m^0$, then all these
functions are divisible by $(z-c)$ [because by hypothesis they are
divisible by a prime divisor of $(z-c)$]. Now by (22)

\weq{31}{-S(\tau) \equiv e_1\tau_1 + e_2\tau_2 + \cdots + e_m\tau_m,}

and hence

\[ -S(\omega) \equiv e_1\tau_1^0 + e_2\tau_2^0 + \cdots + e_m\tau_m^0 \equiv 0 \quad (\mathrm{mod}\ z-c), \]

as was to be proved.

\begin{center}
§ 185. \textbf{\gesperrt{The Fractional Functions in $\Omega$ and the
Taylor Development.}}
\end{center}

A fractional function $\eta$ in the field $\Omega$ can be transformed in
infinitely many ways into an integral function $\mu$ by multiplication
with an integral function $\nu$. We then have
%CONT>>
